% Глава: 17-debugging
% Статус: наполнено (первый прогон).
\chapter{Отладка: gdb и Valgrind}
\label{ch:17-debugging}

\section{Два уровня отладки}

Когда программа «ведёт себя странно», есть два уровня:
\begin{itemize}
  \item \textbf{Логическая ошибка} --- код делает не то,
    что задумано; ищут чтением, принтами, stepping
    в отладчике;
  \item \textbf{Некорректный доступ к памяти} --- код
    обращается к чужой/свободной памяти; ищут
    специализированными инструментами (Valgrind).
\end{itemize}

Отладку ведём на своей ВМ: утилиты \cmd{gdb} и
\cmd{valgrind} работают на сервере \fpath{pglab.nsu.ru}, видя
процесс там, где он запущен.

Для обоих важно собирать с \cmd{-g} (отладочная
информация) и без агрессивной оптимизации
(\cmd{-O0}), чтобы номера строк и значения совпадали
с исходником.

\section{gdb: пошаговая отладка}

\cmd{gdb} --- стандартный отладчик GNU. Типовой цикл:
\begin{minted}[bgcolor=gray!6]{bash}
gcc -Wall -Wextra -g -o app app.c
gdb ./app
\end{minted}

Основные команды (внутри gdb):
\begin{center}
\begin{tabular}{lp{0.6\linewidth}}
\toprule
\textbf{Команда} & \textbf{Действие} \\
\midrule
\cmd{run}      & запустить программу \\
\cmd{break main} & поставить точку останова в main \\
\cmd{break file.c:42} & точка останова по строке \\
\cmd{next}     & следующий шаг (не входя в функции) \\
\cmd{step}     & шаг с заходом в функции \\
\cmd{continue}& продолжить до следующей точки \\
\cmd{print x}  & вывести значение переменной \\
\cmd{backtrace}& показать стек вызовов \\
\cmd{quit}     & выйти из gdb \\
\bottomrule
\end{tabular}
\end{center}

Полезно: \cmd{info locals} --- локальные переменные,
\cmd{watch var} --- остановка при изменении переменной.

\section{Valgrind: поиск проблем с памятью}

\cmd{valgrind} --- инструмент статического анализа
памяти во время выполнения. \cmd{memcheck} ---
наиболее полезный режим:
\begin{minted}[bgcolor=gray!6]{bash}
valgrind --leak-check=full ./app
\end{minted}

Отчёт Valgrind показывает:
\begin{itemize}
  \item \textbf{definitely lost} --- точно утекла память
    (нет ни одного указателя на неё);
  \item \textbf{indirectly lost} --- недоступна через
    другую утечку;
  \item \textbf{still reachable} --- блок жив, но нет
    указателей (обычно нормально на выходе);
  \item \textbf{invalid read/write} --- обращение за
    границами выделенного блока.
\end{itemize}

\begin{tipbox}
Запускайте Valgrind на \emph{отладочной} сборке
(\cmd{-g -O0}): он сам отключает оптимизации, но
с вашими флагами отчёт будет точнее.
\end{tipbox}

\section{Типовые ошибки памяти}

\begin{center}
\begin{tabular}{lp{0.6\linewidth}}
\toprule
\textbf{Ошибка} & \textbf{Что это} \\
\midrule
Buffer overflow   & запись за конец массива/строки \\
Use-after-free    & обращение к памяти после free \\
Dangling pointer  & указатель на уничтоженный объект \\
Double free       & дважды вызвать free на один блок \\
Off-by-one        & индекс на единицу не так (0..n вместо 0..n-1) \\
\bottomrule
\end{tabular}
\end{center}

Все эти ошибки в C «не падают немедленно»: программа
может работать нормально, а потом дать странный
результат или упасть. Valgrind находит их
детерминированно.

\section{Процесс отладки}

\begin{figure}[htb]
  \centering
  \begin{tikzpicture}[node distance=1.4cm]
    \node[data] (bug) {баг};
    \node[process, below=of bug] (repro) {воспроизвести};
    \node[process, below=of repro] (isolate) {изоляция (бинарный поиск, принты)};
    \node[process, below=of isolate] (gdb) {gdb / valgrind};
    \node[process, below=of gdb] (fix) {фикс + тест};
    \draw[flow] (bug) -- (repro);
    \draw[flow] (repro) -- (isolate);
    \draw[flow] (isolate) -- (gdb);
    \draw[flow] (gdb) -- (fix);
  \end{tikzpicture}
  \caption{Типовой цикл отладки.}
  \label{fig:debug-cycle}
\end{figure}

\begin{enumerate}
  \item \textbf{Воспроизвести} стабильно --- без
    воспроизведения отладка невозможна.
  \item \textbf{Изолировать}: бинарный поиск по коду
    (убрать половину), \cmd{printf} или \cmd{assert} в
    подозрительных местах.
  \item \textbf{Наблюдать} через gdb (значения) и
    Valgrind (память).
  \item \textbf{Починить} и \textbf{написать тест},
    чтобы не повторилось (\nameref{ch:15-testing-optimization}).
\end{enumerate}

\section{Типичные ошибки в самом отладчике}
\begin{itemize}
  \item Запустили Valgrind на оптимизированной сборке ---
    отчёт «гуляет»; пересоберите с \cmd{-O0 -g}.
  \item Не поставили \cmd{-g} --- в gdb видно только
    адреса, а не строки.
  \item Ищут баг, который невозможно воспроизвести ---
    сначала стабилизировать условия (фиксированные
    данные, seed для генератора случайных чисел).
\end{itemize}

\section{Связь с остальным курсом}

Отладка опирается на понимание памяти
(\nameref{ch:16-virtual-memory}) и процессов
(\nameref{ch:12-processes}). Valgrind в Docker
(\nameref{ch:19-docker-basics}) требует
\fpath{/proc} и работает медленнее, но применим;
тесты после фикса --- \nameref{ch:15-testing-optimization}.

\section{Задачи для самопроверки}
\begin{enumerate}
  \item Какие флаги \cmd{gcc} нужны для отладочной
    сборки?
  \item Чем \cmd{step} отличается от \cmd{next} в gdb?
  \item Что такое «definitely lost» в отчёте Valgrind?
  \item Приведите пример off-by-one ошибки в C.
  \item Почему Valgrind медленнее обычного запуска и
    когда это допустимо?
\end{enumerate}
